\section{Forget-Set Audit}
\label{sec:audit}

If side-effects depend on the forget set, then choosing $\Df$ is itself a tunable design decision --- but \emph{evaluating} that decision currently requires running unlearning. We now ask whether $\Df$'s intrinsic geometry is a sufficient signal to predict the three-layer corruption profile without ever fine-tuning.

\subsection{Audit Problem}
\label{sec:audit-problem}

\paragraph{Inputs.} The $50\!\times\!384$ embedding matrix of a candidate $\Df$, summarised by the $12$-dim geometric descriptor of Sec.~\ref{sec:features}. \emph{No} unlearned checkpoint, \emph{no} evaluation text, \emph{no} retain set is required.

\paragraph{Outputs.} Three scalars: predicted geo-mean $\log r$ for $\Lone$, $\Ltwo$, $\Lthree$. We treat audit primarily as a \emph{ranking} problem (which forget sets to flag), and report regression metrics for completeness.

\paragraph{Predictor.} Ridge regression with $\alpha{=}1.0$, one regressor per layer, evaluated under leave-one-cluster-out (LOO over the $10$ HDBSCAN clusters; no triplet sees its cluster during fitting). The baseline predicts the mean of the $9$ training-cluster targets.

\subsection{Audit Performance}
\label{sec:audit-perf}

Table~\ref{tab:audit} reports held-out $R^2$, MAE, and Spearman $\rho$, plus $95\%$ percentile bootstrap CIs for $\rho$ ($n_{\text{boot}}{=}10{,}000$, seed $0$, bootstrapping over the $10$ LOO (true, predicted) cluster pairs).

\begin{table}[t]
\centering
\small
\setlength{\tabcolsep}{3.5pt}
\begin{tabular}{@{}lccc@{}}
    \toprule
                                  & $\Lone$            & $\Ltwo$            & $\Lthree$            \\
    \midrule
    Audit $R^2$                   & $\mathbf{+0.246}$  & $\mathbf{+0.705}$  & $\mathbf{+0.215}$    \\
    Baseline $R^2$ (LOO mean)     & $-0.020$           & $-0.020$           & $-0.020$             \\
    Audit MAE                     & $\mathbf{0.141}$   & $\mathbf{0.050}$   & $\mathbf{0.020}$     \\
    Baseline MAE                  & $0.169$            & $0.105$            & $0.022$              \\
    \midrule
    Spearman $\rho$               & $+0.531$           & $+0.840$           & $+0.491$             \\
    $\rho$ 95\% CI                & $[+0.37, +0.67]$   & $[+0.76, +0.89]$   & $[+0.33, +0.63]$     \\
    \bottomrule
\end{tabular}
\caption{Forget-set audit performance. LOO over $10$ clusters; Ridge ($\alpha{=}1$) on $12$-dim forget-set geometry. All three $\rho$ CIs exclude zero; $\Ltwo$ is the strongest layer.}
\label{tab:audit}
\end{table}

\begin{figure}[t]
\centering
\includegraphics[width=0.95\columnwidth]{fig_audit_scatter.pdf}
\caption{\textbf{Audit scatter}, predicted vs.\ true geo-mean $\log r$. Each point is one held-out forget set under LOO. $\Ltwo$ is closest to the diagonal ($\rho{=}0.84$); $\Lone$ and $\Lthree$ are noisier but rank-significant.}
\label{fig:audit-scatter}
\end{figure}

\begin{figure}[t]
\centering
\includegraphics[width=0.95\columnwidth]{fig_forget_spread.pdf}
\caption{\textbf{Forget-set spread.} Distribution of per-forget-set $\Lone$ across the $100$ triplets, grouped by HDBSCAN cluster. The variation auditing must capture is between-cluster more than within-cluster.}
\label{fig:forget-spread}
\end{figure}

\textbf{$\Ltwo$ is the strongest layer.} With $R^2{=}0.71$ and $\rho{=}0.84$, the auditor recovers $71\%$ of locality variance from forget-set geometry alone --- consistent with the intuition that locality damage tracks how concentrated and self-similar $\Df$ is in embedding space.

\textbf{$\Lthree$ is non-trivial.} On a coarser problem ($n{=}10$ in early experiments) the L3 audit had $R^2{=}-0.46$; at $n{=}100$ the same predictor reaches $R^2{=}+0.22$ with $\rho{=}0.49$ ($95\%$ CI exclude zero). Statistical power, not signal, was the previous bottleneck.

\textbf{Audit cost.} Computing the $12$-dim feature for one $\Df$ takes milliseconds; fitting the Ridge model on $99$ training triplets is sub-second; predicting on a new $\Df$ is essentially instant. By contrast, evaluating one candidate $\Df$ via the standard pipeline costs one unlearn run ($\approx 75$\,s on H100) plus one column of cross-PPL ($\approx 220$\,s for $100$ eval pairs). At $n{=}100$ candidates the full grid is $\approx 7$ GPU-hours; the auditor runs in seconds on a CPU.

\subsection{Positioning: a Pre-Screen, Not a Replacement}
\label{sec:positioning}

We position the auditor as a \emph{warning tool} for the forget-set selection step, not as a substitute for the actual unlearning evaluation.

\begin{table}[t]
\centering
\small
\setlength{\tabcolsep}{3pt}
\renewcommand{\arraystretch}{1.05}
\begin{tabular}{@{}p{0.18\columnwidth} p{0.36\columnwidth} p{0.36\columnwidth}@{}}
    \toprule
              & \textbf{Can}                                  & \textbf{Cannot} \\
    \midrule
    Rank      & flag ``which $\Df$ is worse'' ($\rho$ CIs $\!>\!0$, $\Ltwo$ $0.84$) & assign trustworthy absolute PPL values \\
    Screen    & red/yellow/green warnings on candidates       & score an unlearning algorithm \\
    Compute   & seconds vs.\ several GPU-hours                & replace the final unlearn run \\
    \bottomrule
\end{tabular}
\caption{Audit positioning: cheap pre-screen for forget-set selection.}
\label{tab:positioning}
\end{table}

\paragraph{Workflow.} Given $N$ candidate forget sets, audit all $N$ in seconds, sort by predicted layer profile, and run real unlearning only on the top-$k$. At $N{=}100$ and $k{=}10$, the pre-screen reduces the unlearning budget by $10\times$ while retaining the worst-case offenders identified by the auditor.
